\section{Data}
\label{sec:data}

We assemble a panel covering the universe of Brazilian municipalities,
their public hospital network, and an extended window of inpatient
admissions ($2010$ through $2024$, $15$ years) and mortality outcomes
($2010$ through $2023$, $14$ years; the latter ends in $2023$, the last fully consolidated SIM vintage at the time of
writing). All sources are administrative microdata maintained by the
Brazilian Ministry of Health (DATASUS) and the IBGE statistical office.

\subsection{Hospital Admissions: SIH-AIH}

The Sistema de Informações Hospitalares do SUS (SIH-AIH) records every
inpatient admission billed to the federal SUS reimbursement system.
Each Autorização de Internação Hospitalar (AIH) corresponds to one
hospital stay. Unique linkage uses the AIH-1 (principal hospitalization)
record; AIH-5 (continuation) records are excluded to avoid
double-counting. Variables include the patient's $6$-digit
municipality of residence (\texttt{MUNIC\_RES}), the establishment's
CNES identifier, the principal ICD-10 diagnosis (\texttt{DIAG\_PRINC}),
the realized procedure code (\texttt{PROC\_REA}), in-hospital death
flag (\texttt{MORTE}), age and sex, dates of admission and discharge,
and the federally-paid amount.

We retain $179{,}499{,}614$ admissions with non-missing residence
municipality, hospital CNES, and competence year over $2010\text{--}2024$.
We restrict to admissions in establishments classified as hospitals in
the CNES master file (excluding non-hospital types such as primary care
units, clinics, and specialty offices), retaining
$179{,}420{,}885$ admissions ($\approx 100\%$ of the gross). A small
fraction is lost to ineligible non-hospital codes.

For graph construction (Section \ref{sec:method}), we aggregate the
admission-level data to the municipality--hospital--year level, recording
the count of admissions, the count of in-hospital deaths, and the
average length of stay and reimbursement value. The resulting bipartite
edge list contains $2{,}703{,}889$ municipality-hospital-year cells; the
period-pooled version aggregating $2010\text{--}2024$ contains
$648{,}973$ unique municipality-hospital edges.

\subsection{Mortality: SIM}

The Sistema de Informação sobre Mortalidade (SIM) records every death
certificate issued in Brazil. Variables include the underlying ICD-10
cause of death (\texttt{CAUSABAS}), municipality of residence
(\texttt{CODMUNRES}), municipality of occurrence (\texttt{CODMUNOCO}),
age (encoded with a leading hint character distinguishing days, months,
and years), sex, race, and educational attainment of the mother (for
infant deaths). Coverage of total mortality is estimated by the
Ministry of Health to be $\approx 95\text{--}99\%$ in the South and
Southeast and $\approx 85\text{--}95\%$ in the North; ill-defined causes
(R-codes) range from $\approx 5\%$ in the South to $\approx 15\%$ in
the rural Northeast.

\paragraph{Amenable mortality.}
We construct the rate of amenable mortality at the municipality-year
level using the consensus list of \citet{nolteMckee2008}, updated by
the OECD \citep{oecd2019}, of conditions for which death below age $75$
is considered avoidable through timely and effective medical care.
The list comprises $\approx 60$ ICD-10 categories spanning
ischemic heart disease, stroke, sepsis, perforated appendix, treatable
cancers, complications of childbirth, and other conditions. We follow
the standard convention of restricting to deaths under age $75$ and
expressing the rate per $100{,}000$ population.

Over $2010\text{--}2023$, the SIM records $18{,}878{,}605$ deaths total,
of which $5{,}362{,}418$ ($28.4\%$) are amenable under our definition.
Population denominators come from IBGE annual estimates (Tabela 6579 of
SIDRA for $2010\text{--}2021$ and the $2022$ Census plus
$2023\text{--}2024$ estimates from \texttt{estimativa\_dou}); for years
without SIDRA estimates we use the IBGE annual interpolations.

\subsection{Hospital Master and Habilitations: CNES}

The Cadastro Nacional de Estabelecimentos de Saúde (CNES) records the
universe of healthcare facilities at monthly granularity. We construct
two derived tables.

The first is a hospital master file aggregating CNES record
characteristics across $2005\text{--}2025$, retaining for each CNES the
modal municipality of operation, modal facility type
(\texttt{CO\_TIPO\_ESTABELECIMENTO}), modal management sphere (federal,
state, municipal, or private contract), and the first and last
competence months observed. We retain establishments classified as
hospitals: general hospital (\texttt{05}), specialty hospital
(\texttt{07}), and day hospital (\texttt{62}). The resulting universe
contains $15{,}316$ hospital CNES.

The second is a habilitation panel constructed from the CNES habilitation
files (\texttt{HB} layer). Each row records a hospital--habilitation--
month observation. We classify habilitations into the seven
high-complexity categories listed in Section \ref{sec:setting} using a
lookup table from the SIGTAP procedure-group codes (\texttt{SGRUPHAB})
to our category labels. A hospital is classified as
\emph{high-complexity active} in our panel if it holds at least one
high-complexity habilitation that is active in any month during
$2010\text{--}2024$. This identifies $10{,}690$ hospitals
($69.8\%$ of the master universe).

\subsection{Spatial and Auxiliary Data}

Municipality centroids and shapefiles come from the IBGE
\texttt{municipios\_2010} layer ($5{,}565$ municipalities; we use the
$6$-digit IBGE code for SIH/SIM compatibility). State shapefiles come
from \texttt{estados\_2010}. Both layers are projected on SIRGAS 2000
(EPSG:4674) for distance computations and reprojected to a Lambert
azimuthal equal-area for visualization.

GDP at the municipality level is from IBGE Tabela $5938$ (Produto
Interno Bruto dos Municípios) for $2010\text{--}2021$ in current
nominal reais, deflated to $2020$ R\$ using the IPCA. Population is
from Tabela $6579$ for $2010\text{--}2021$ and the $2022$ Census plus
official annual estimates for $2023\text{--}2024$.

\subsection{Final Sample}

After matching, $5{,}588$ municipalities of residence appear in
SIH-AIH $2010\text{--}2024$. Joining to the IBGE $2010$ shapefile
($5{,}565$ municipalities) yields $5{,}565$ municipalities with valid
centroids, non-missing population and GDP, and at least one neighbor
in the patient-flow graph. This is the analysis sample for the
divergence map and the descriptive results; the $23$ residual
municipalities (recently created codes not in the $2010$ IBGE layer
or with incomplete digital records) drop out. Of the $15{,}316$ hospital CNES,
$6{,}710$ ($43.8\%$) participate in the bipartite graph after the SIH
universe filter, and $3{,}488$ ($32.6\%$ of the $10{,}686$
high-complexity hospitals) hold a high-complexity habilitation
\emph{and} appear in the embedding. The hospital coverage gap is one of the limitations
discussed in Section \ref{sec:discussion}: hospitals with active
SIGTAP habilitations but no SUS-billed admissions during the panel
window---typically philanthropic and small private hospitals serving
non-SUS clients---are not embedded and therefore are excluded from the
network metric.
